% ============================================================================
%  model_cory: specification
%  enightmodels::model_legacy(), as called with model_legacy_config() defaults.
%  Source: enightmodels:::model_legacy_total and :::model_legacy_cand.
%  Compile: pdflatex cory_model_spec.tex
% ============================================================================
\documentclass[11pt]{article}

\usepackage[margin=1in]{geometry}
\usepackage{amsmath, amssymb}
\usepackage{microtype}
\usepackage{enumitem}
\usepackage[colorlinks=true, allcolors=blue]{hyperref}

\setlist[itemize]{itemsep=2pt, topsep=4pt}
\newcommand{\var}[1]{\texttt{#1}}
\newcommand{\clamp}{\operatorname{clamp}}

\title{\texttt{model\_cory}: specification}
\author{Young-Kyung Kim}
\date{\today}

\begin{document}
\maketitle

\noindent
This documents \texttt{enightmodels::model\_legacy()} as it is called by
\texttt{scripts/model/live/R/model\_cory.R}, with \texttt{model\_legacy\_config()} defaults. It is the production legacy model
and the second of the two baselines the precinct $\times$ vote-mode model is compared
against. Sections follow the same order as the main model's specification. The source is
\texttt{enightmodels:::model\_legacy\_total} and
\texttt{enightmodels:::model\_legacy\_cand}; this document describes those functions, not
a reimplementation of them.

\paragraph{Configuration.} The \texttt{model\_legacy\_config()} defaults in force:

\begin{center}
\small
\begin{tabular}{@{}ll@{\hspace{2.5em}}ll@{}}
\var{by\_method} & \var{TRUE}        & \var{min\_unit\_rep} & $0.05$ \\
\var{strata}     & \var{geostratum}  & \var{tr\_pre}        & $1$    \\
\var{stage\_complete} & $6$          & \var{tr\_min}        & $0.25$ \\
\var{min\_complete\_per\_stratum} & $1$ & \var{tr\_max}     & $4$    \\
\end{tabular}
\end{center}

\section{Outcomes}

The model runs in two stages. The first predicts \textbf{total votes} per
geo $\times$ mode; the second predicts \textbf{each candidate's votes} using the first
stage's reporting fractions. Both are ratio estimators against a past-election baseline
$\var{past}$, computed within a stratum. The stratum is the \var{geostratum}, not the
county.

Unlike the other two models, nothing is estimated on a share scale and nothing is
modeled as a share: the model predicts candidate \emph{vote counts} directly.

\section{Turnout: total votes}

\paragraph{Global pre-scale.} Before anything else, the baseline is rescaled so that its
statewide total matches what has been counted so far, floored at one:
\begin{equation}
  t = \max\!\left(\frac{\sum_{\text{all cells}} \text{votes}}{\sum_{\text{all cells}} \var{past}},\ 1\right),
  \qquad
  \var{past} \leftarrow \var{past} \cdot t .
  \label{eq:prescale}
\end{equation}
The floor at $1$ means the baseline is never scaled \emph{down}. Early in the night the
counted total is small, so $t = 1$ and the step is inert; it becomes active only once
counting is far enough along that the running total exceeds the baseline.

\paragraph{Completeness.} Cells are aggregated to
geo $\times$ mode $\times$ stratum, and a unit is complete when
$\var{stage} \geq 6$. This is a reported flag, not something inferred from the counts.

\paragraph{Stratum inflation.} Strata containing at least one complete unit get a single
inflation factor per mode,
\begin{equation}
  \begin{aligned}
  \var{infl}_{\sigma m}
    &= \clamp\!\left(
      \frac{\sum_{i \in \sigma} \var{past}_{im} \cdot (\text{votes}_{im} / \var{past}_{im})}
           {\sum_{i \in \sigma} \var{past}_{im}},\ 0.25,\ 4
    \right) \\[4pt]
    &= \clamp\!\left(
      \frac{\sum_{i \in \sigma} \text{votes}_{im}}{\sum_{i \in \sigma} \var{past}_{im}},\ 0.25,\ 4
    \right),
  \end{aligned}
  \label{eq:infl}
\end{equation}
a $\var{past}$-weighted mean of the votes-to-baseline ratio. Strata with no complete unit
take $\var{infl} = 1$. Prediction then substitutes the count where a unit is complete and
inflates the baseline where it is not:
\begin{equation}
  \widehat{T}_{im} =
  \begin{cases}
    \text{votes}_{im} & \text{complete,}\\
    \var{past}_{im} \cdot \var{infl}_{\sigma(i) m} & \text{otherwise.}
  \end{cases}
  \label{eq:predtot}
\end{equation}

\paragraph{Excess redistribution.} A final per-geo pass reconciles predictions against
counts that have overtaken them. With $x_{im} = \max(\text{votes}_{im} - \widehat{T}_{im},\,0)$
the overshoot in a cell, given $5\%$ headroom when the unit is not complete,
\begin{equation}
  \widehat{T}_{im} \leftarrow \max\!\left(\frac{\sum_{m'} \text{votes}_{im'}}{\sum_{m'} \widehat{T}_{im'}},\ 1\right) \widehat{T}_{im},
  \qquad
  \widehat{T}_{im} \leftarrow \widehat{T}_{im} +
  \begin{cases}
    x_{im} & x_{im} > 0,\\[3pt]
    -\dfrac{\widehat{T}_{im}}{\sum_{x=0} \widehat{T}} \displaystyle\sum_{m'} x_{im'} & x_{im} = 0 .
  \end{cases}
  \label{eq:excess}
\end{equation}
Cells that have overshot are raised to their counted value; the cost is taken
proportionally from the cells that have not. The geo total is conserved by construction.
This is the model's one adding-up guarantee, and it operates on overshoot only --- there
is no corresponding treatment of cells that are short.

\section{Candidates}

The second stage uses the first stage's reporting fraction
$\rho_{im} = \text{votes}_{im} / \widehat{T}_{im}$ to scale each candidate's past votes.
Within a stratum $\times$ mode $\times$ candidate, over units that have reported at least
$\var{min\_unit\_rep} = 5\%$,
\begin{equation}
  \var{ratio}_{\sigma m k} =
  \frac{\sum_{i} \text{votes}_{imk} / \rho_{im}}{\sum_{i} \var{past}_{imk}},
  \qquad
  \widehat{v}_{imk} =
  \begin{cases}
    \text{votes}_{imk} & \var{stage} \geq 6,\\
    \var{past}_{imk} \cdot \var{ratio}_{\sigma m k} & \text{otherwise,}
  \end{cases}
  \label{eq:cand}
\end{equation}
with the fallback $\var{ratio} = \widehat{T}_{im} / \sum_k \var{past}_{imk}$ where a
stratum has no qualifying unit. Dividing by $\rho_{im}$ grosses each reporting unit up to
its implied final count before the ratio is formed, so partially-counted units contribute
on the same footing as complete ones.

\section{Uncertainty}

\textbf{There is none.} \texttt{model\_legacy()} returns \var{pred\_min} and
\var{pred\_max} as \texttt{NA}, and no draws, interval or variance are produced anywhere
in either stage. Every \texttt{model\_cory} interval column is consequently \texttt{NA}, and the model contributes a point estimate only.

\section*{Limitations}
\addcontentsline{toc}{section}{Limitations}

Stated plainly. Several of these are deliberate: this is a long-serving production model
whose simplicity and predictability are part of its value.

\begin{itemize}
  \item \textbf{No uncertainty at all.} The model cannot be assessed on coverage or
        interval width, and cannot contribute a probability of winning. Any comparison
        against it is necessarily a point-accuracy comparison.
  \item \textbf{No ceiling.} Nothing bounds predicted votes by registrants. The clamp in
        \eqref{eq:infl} bounds the \emph{ratio} at $4\times$ the baseline, which is a
        loose constraint and not a registrant-based one.
  \item \textbf{Candidate predictions need not sum to the turnout prediction.} The two
        stages are separate ratio estimators, and \eqref{eq:cand} applies a per-candidate
        factor with no constraint tying $\sum_k \widehat{v}_{imk}$ to $\widehat{T}_{im}$.
  \item \textbf{The inflation factor is computed over all units in a qualifying
        stratum.} A stratum qualifies on having one complete unit, but \eqref{eq:infl}
        then averages over every unit in it, including barely-started ones whose
        votes-to-baseline ratio is near zero. This biases $\var{infl}$ \emph{downward}
        early in the night, when few units are complete --- the opposite direction to the
        usual partial-count failure.
  \item \textbf{Clamping rather than flagging.} $\clamp(\cdot, 0.25, 4)$ in
        \eqref{eq:infl} silently transforms an out-of-range ratio. A ratio outside
        $[0.25, 4]$ is a strong signal that the baseline is wrong for that stratum, or
        that the feed is misaligned; clamping discards that signal rather than surfacing
        it.
  \item \textbf{Completeness is load-bearing.} Both stages branch on
        $\var{stage} \geq 6$, so the model depends on a feed-supplied flag that in
        production arrives late, uses different vocabulary per state, and can move
        backwards. There is no fallback that degrades gracefully if the flag is wrong.
  \item \textbf{The stratum is the geostratum, not the county.} Dependence between
        precincts in the same county is not addressed, and no shrinkage is applied to
        stratum effects: \eqref{eq:infl} is an unpenalized stratum mean.
\end{itemize}

\end{document}
